\documentclass{article}
\usepackage{amsmath}
\begin{document}
\section*{Two evidence sources can discharge one gate duty, not renew it}
Keep the native pose stall and the locally qualified optical tube.
The courier's actual physical witness always queues a packet, and
the independent diagnostic always queues another. At the gate,
either its own qualified fiducial observation or the native delivery
of the courier's whole-body clearance receipt can authorize the same
single timed closure. The first usable evidence is a source event,
not a prediction that a receipt will be delivered. The optical
branch is chosen only when the qualified modeled local observation
is actually available. Without it, no gate effect is installed
before a valid packet. Source binding, active authority, single
mission identity and monotone non-reentry movement are provisions.

With live local sensing and delayed broker repair, the local
clearance is used before the pending packets deliver. When the
courier's second authorization actually reaches the gate on
repair, it resolves to the same immutable SQLite effect result.
The exact stored effect interval, decision rows and durable event
counts are unchanged. Crashing before or after the first effect
commit does not turn the native duplicate into a second closure
or renew the original duration. Both the separate diagnostic and
the exact later dock result still complete. The control retains
the extra radio work even when it is redundant; this is not a
bandwidth gain over knowing in advance that local sensing works.
The actual courier pose, gate and allowed breadcrumb points all
lie on the outside of the installed slab when the pending packets
dispatch; no general dynamic radio model is implied.

With a modeled optical blackout and the same physical path, faster
radio repair supplies the only qualifying observation in time.
Slower repair still yields correct gate authorization but misses
the separate timed-work cutoff. All 24 variants in each arm
preserve egress, diagnostic and exact dock work; only the two
timely evidence scenarios satisfy the timed closure in all 24.
The result is a conditional avoidance of one evidence bottleneck
and duplicate actuator effect, not dominance over a direct
first-qualified-source controller with the same authority.
This remains a held-in ordinary-program replay, not independent
mission acceptance or a claim about deployed sensors.
\end{document}
